\section{Introduction}
\label{sec:intro}

\subsection{Project context}
Scene understanding for autonomous driving is commonly cast as a dense
labelling problem. \emph{Semantic segmentation} assigns a class label to
every pixel; \emph{instance segmentation} additionally separates individual
objects of the same class; and \emph{panoptic segmentation}~\cite{kirillov}
unifies the two by labelling every pixel with both a class and, for
countable ``thing'' classes, an instance identity. Deep networks have made
these tasks remarkably accurate in structured, closed-world conditions where
the set of categories is fixed in advance.

Real driving, however, is an \emph{open world}. Networks trained on a closed
label set are by construction unable to represent objects outside that
set, yet such unknown, anomalous or out-of-distribution (OoD) objects---a
lost cargo box, an animal, debris on the road---are exactly the cases where
a perception failure is most dangerous~\cite{smiyc,fishyscapes}. Detecting
and localising these regions is therefore a safety-critical complement to
ordinary segmentation. This project follows a single thread through both
problems: it studies the dominant \emph{mask-classification} family of
segmentation models, compares two pre-trained variants of a state-of-the-art
member of that family, adapts one of them to the driving domain by
fine-tuning, and finally repurposes the same models for anomaly segmentation
using training-free (post-hoc) scoring rules.

\subsection{Literature foundations}
\paragraph{From per-pixel to mask classification.}
Efficient per-pixel networks such as ERFNet~\cite{erfnet} established that
accurate real-time semantic segmentation is feasible on embedded automotive
hardware by factorising convolutions for a favourable accuracy/latency
trade-off. Such models treat segmentation as independent per-pixel
classification. MaskFormer~\cite{maskformer} challenged this view: it shows
that predicting a set of binary \emph{masks}, each paired with a single class
label, is general enough to solve semantic, instance and panoptic
segmentation with one model, loss and training procedure.
Mask2Former~\cite{mask2former} refined this idea with \emph{masked
attention}, restricting cross-attention to the region of each predicted mask,
which improves both accuracy and convergence and made a single architecture
competitive across all three tasks.

\paragraph{Foundation backbones and EoMT.}
In parallel, self-supervised pre-training produced strong general-purpose
visual features: DINOv2~\cite{dinov2} learns robust representations from
unlabelled images that transfer across many dense tasks. The Encoder-only
Mask Transformer (EoMT)~\cite{eomt} brings these threads together: it argues
that a sufficiently large, well pre-trained plain Vision Transformer does not
need the task-specific adapters, pixel decoders and Transformer decoders of
Mask2Former. Instead, learnable queries are processed \emph{inside} the ViT
encoder itself, yielding a simpler architecture that matches specialised
models while being substantially faster. EoMT is the central model of this
project, instantiated with a DINOv2 ViT backbone.

\paragraph{Anomaly segmentation and post-hoc scoring.}
The anomaly-segmentation task and its evaluation are defined by the
SegmentMeIfYouCan (SMIYC)~\cite{smiyc} and Fishyscapes~\cite{fishyscapes}
benchmarks, together with the earlier RoadAnomaly dataset~\cite{roadanomaly};
all judge whether a pixel belongs to the known Cityscapes~\cite{cityscapes}
classes or is anomalous. A practical and popular strategy is to derive an
anomaly score \emph{post-hoc} from an off-the-shelf segmentation network,
with no retraining. Classic rules are Maximum Softmax Probability
(MSP)~\cite{msp} and Max-Logit / energy-style scores from the scaling study
of Hendrycks~\etal~\cite{maxlogit}. For mask-classification models a tailored
rule, Rejected-by-All (RbA)~\cite{rba}, scores a pixel as anomalous when
\emph{no} object query claims it with confidence---a natural fit for
set-prediction architectures. Finally, when adapting large models under tight
compute budgets, parameter-efficient methods such as LoRA~\cite{lora} are an
attractive alternative to full fine-tuning.

\paragraph{Contributions.}
Within this context, the report (i) sets up a fair, shared semantic protocol
to compare a Cityscapes-trained and a COCO-trained~\cite{coco} EoMT despite
their different label spaces (\cref{sec:method,sec:results}); (ii) fine-tunes
the COCO model on Cityscapes with a staged freezing schedule and analyses the
resulting mIoU; and (iii) benchmarks pixel-based (ERFNet) and mask-based
(EoMT) post-hoc anomaly detection with MSP, Max-Logit, Max-Entropy and RbA,
including temperature scaling.
